\section{Discussion}
\label{sec:discussion}

\para{What do the results mean?} The experiment supports the main hypothesis, but with a stronger conclusion than we expected. In this setup, benign-refusal fine-tuning does not just increase refusal on semantically related prompts. It nearly collapses the model into an always-refuse policy. The 1-shot adapter is already enough to make almost every benign prompt look unsafe, and the 5-shot adapter completes the collapse.

\para{Why does the helpful control matter?} The helpful control is important because it shows that the intervention is not purely label-specific. Training on only five benign examples with ordinary helpful answers still raises benign refusal by 15.2 points and harmful refusal by 9.2 points. This suggests that tiny-data SFT on a small aligned model can itself destabilize the refusal boundary. Even so, the refusal-labeled conditions are far stronger, so the user's intended intervention clearly amplifies over-refusal well beyond the generic fine-tuning effect.

\para{Connection to prior work.} Our findings fit naturally with both latent-mechanism and safety-fragility accounts. If refusal is partly mediated by a shared direction \cite{arditi2024refusal}, then a small amount of concentrated refusal supervision could move the model along that direction for many unrelated prompts. If aligned models are broadly brittle under narrow post-training updates \cite{qi2023finetuning,betley2025emergent}, then the near-global shift we observe is less surprising. What our experiment adds is a direct causal demonstration using an extremely small benign-refusal dataset.

\para{Broader implications.} The practical implication is that downstream safety behavior can be fragile in both directions. Much prior work emphasizes under-refusal after fine-tuning. Our results show the complementary risk: a tiny benign-refusal edit can make an aligned assistant broadly unusable. This matters for adapter release, narrow enterprise tuning, and safety patching workflows that relabel small examples without broad regression testing.

\para{Limitations.} This study has several clear limits. We use a small 1.5B model, which may be more brittle than larger production systems. The training regime is intentionally extreme relative to the dataset size: 100 steps on 1 or 5 examples can saturate behavior quickly. The helpful control also shifts refusal, so not all of the effect is uniquely attributable to refusal labels. Our primary labeler is regex-based, and the API judge validates only a subset of outputs. In addition, \phtest and \orbench are sampled rather than evaluated in full, and we do not run a multi-seed sweep.

\para{What should come next?} The most important follow-up is to map the boundary of the effect. We need larger models, more seeds, and finer sweeps over step count and shot count to tell whether the observed collapse is a general property of aligned models or a particularly brittle small-model regime. It is also worth testing mitigation methods such as refusal reflection or vector ablation after the induced over-refusal appears.
